% ===================================================================
%  QUADRO N.º 7 — A aldeia no inverno. — A casa de lavoura na primavera.
%
%  Ceci n'est PAS une transcription : c'est une traduction, faite en
%  2026, en portugais d'aujourd'hui. Voir texto/pt/10-quadro-01.tex
%  pour la consigne, qui vaut pour les seize fichiers.
% ===================================================================

%%K t07-noto-1 noto
\VUnotes{143.2mm}{%
(*) Para o pormenor de uma refeição, ver os quadros 6 e 13.%
}

%%K t07-apar-1 apar
\VUsaut{4.0mm}
\VUcentre{13.2pt}{90}{SEGUNDA SÉRIE}
\VUsaut{3.0mm}
\VUfilet{16mm}
\VUsaut{5.6mm}
\VUtitre{15.0pt}{60}{QUADRO N.º 7}
\VUsaut{3.0mm}

%%K t07-tit-1 sub
\VUcentre{12.0pt}{0}{\textit{Primeira cena.}}
\VUsaut{3.0mm}

%%K t07-tit-2 sub
\VUpk{10.2pt}{40}{A aldeia no inverno.}
\VUsaut{4.0mm}

%%K t07-c1-01-1 p
1. --- Durante as férias do Ano Novo acompanhei o meu pai a Vilanova, uma
aldeiazinha onde tinha alguns assuntos a tratar.

%%K t07-c1-02-1 p
2. --- Tomámos a \VUgras{diligência} \textsuperscript{(11)} que faz o serviço
entre a cidade e essa aldeia; mas como estava muito frio, aquela viagem
num veículo tão pouco cómodo não foi nada agradável.

%%K t07-c1-03-1 p
3. --- Por isso sentimos verdadeiro prazer quando vimos que o
\VUgras{cocheiro} parava a sua carruagem na praça, diante da
\VUgras{estalagem} \textsuperscript{(2)} do \textit{Urso Branco}. O
\VUgras{estalajadeiro} \textsuperscript{(4)} esperava-nos à soleira da sua casa,
fumando tranquilamente o seu \VUgras{cachimbo}.

%%K t07-c1-03-2 p
Convidou-nos a entrar e eu não me fiz rogado para me instalar diante do
fogo de sarmentos que crepitava na lareira. Então ria-me do áspero
\VUgras{vento do norte} que fazia chiar a velha \VUgras{tabuleta}
\textsuperscript{(3)} e levava em negros remoinhos o \VUgras{fumo}
\textsuperscript{(1)} que saía da chaminé.

%%K t07-c1-04-1 p
4. --- Era a hora de comer. Sem esperar mais, fizemos honra à refeição
simples mas saborosa que nos serviram (*).

%%K t07-tit-3 sub
\VUpk{10.2pt}{40}{Algumas visitas.}
\VUsaut{3.0mm}

%%K t07-c1-05-1 p
5. --- Depois de comer acompanhei o meu pai, que é agente de seguros, a
casa dos seus clientes. Primeiro visitámos o \VUgras{ferreiro}
\textsuperscript{(5)}, que mora junto à estalagem. Estava a ferrar um cavalo.
Admirei a força do seu ajudante, que segurava firmemente contra o seu
avental de couro o \VUgras{casco} do cavalo enquanto o ferreiro fixava a
\VUgras{ferradura} \textsuperscript{(6)} a fortes marteladas.

%%K t07-c1-06-1 p
6. --- Depois fomos a casa do \VUgras{sapateiro} \textsuperscript{(7)} da aldeia, o
tio Jaime. Embora tivesse por tabuleta uma bota \VUgras{magnífica},
contentava-se com pôr solas num par de sapatos velhos e esburacados.

%%K t07-c1-06-2 p
O bom homem disse-nos a rir: «Na cidade era “sapateiro” e não tinha
clientes. Aqui sou “remendão” e há trabalho de sobra».

%%K t07-c1-07-1 p
7. --- Falou-nos do seu vizinho o \VUgras{barbeiro} \textsuperscript{(8)}, cuja
clientela está descontente. As suas \VUgras{navalhas} estão sempre
denteadas e as suas \VUgras{tesouras} nunca cortam. Mas como é o único
barbeiro da aldeia, não há remédio senão deixar-se barbear e cortar o
cabelo por ele. Além disso é um homem avarento e antipático.

%%K t07-c1-08-1 p
8. --- Por sorte, os demais \VUgras{vizinhos} não se parecem com ele. O
\VUgras{carpinteiro de carros} \textsuperscript{(9)}, por exemplo, é um homem excelente,
trabalhador e prestativo. Dá gosto levar-lhe um carro a arranjar. O seu
trabalho está sempre bem acabado.

%%K t07-c1-09-1 p
9. --- O tio Jaime também não se queixava do \VUgras{seleiro}
\textsuperscript{(10)}, a quem às vezes ajuda a coser os \VUgras{arreios} quando
o trabalho aperta.

%%K t07-c1-09-2 p
O meu pai perguntou-lhe pela família. «Estão todos bem. A minha neta está
na \VUgras{escola} \textsuperscript{(12)}. A sua \VUgras{professora}
\textsuperscript{(13)} está muito contente com ela; é boa \VUgras{aluna}
\textsuperscript{(14)}. Não é como aquele malandro do meu filho; só pensa em
brincar. O \VUgras{professor} \textsuperscript{(15)} não consegue ensinar-lhe
nada».

%%K t07-tit-4 sub
\VUsaut{3.0mm}
\VUpk{10.2pt}{40}{Conversas da aldeia.}
\VUsaut{3.0mm}

%%K t07-c1-10-1 p
10. --- O bom homem contou-nos depois as novidades da aldeia. Iam
construir uma escola nova, porque a que estava instalada na \VUgras{casa
da câmara} tinha ficado muito pequena. De resto era fácil, porque bastava
comprar uma parte do jardim do \VUgras{moleiro}. Isso viria muito a
propósito ao pobre homem. Desde que chegaram os frios, as \VUgras{mós} do
\VUgras{moinho} \textsuperscript{(16)} já não podiam moer o trigo, porque a
\VUgras{roda} \textsuperscript{(17)} estava presa no \VUgras{gelo}
\textsuperscript{(25)} do \VUgras{ribeiro} \textsuperscript{(23)}.

%%K t07-c1-11-1 p
11. --- «E por falar em construções, viram a nossa nova \VUgras{igreja}
\textsuperscript{(18)}? Está muito pitoresca sob o seu manto de neve. Os
\VUgras{corvos} \textsuperscript{(19)}, que já não reconhecem o seu velho
campanário, pousam tristemente na grande árvore do \VUgras{cemitério}
\textsuperscript{(20)}».

%%K t07-c1-12-1 p
12. --- A ideia do cemitério lembrou ao sapateiro as pessoas que o meu
pai conhecera e que tinham morrido no ano passado. «Quantas
\VUgras{campas} \textsuperscript{(21)}, meu bom senhor, se juntaram às outras sob
o \VUgras{cipreste} \textsuperscript{(22)}! A morte levou-nos o nosso velho
notário. Toda a aldeia assistiu ao seu \VUgras{enterro}».

%%K t07-c1-13-1 p
13. --- «Ali está o seu sucessor, a patinar no ribeiro. Veio há pouco
para que lhe arranjasse a correia do seu \VUgras{patim} \textsuperscript{(26)},
que se tinha rompido».

%%K t07-c1-14-1 p
14. --- «Ali está também, à beira do ribeiro, o novo \VUgras{guarda rural}
\textsuperscript{(27)}. É um antigo militar da guarda e o terror dos garotos da
aldeia. Fogem assim que veem as suas \VUgras{polainas} \textsuperscript{(29)} e o
seu \VUgras{quépi} \textsuperscript{(28)}».

%%K t07-c1-15-1 p
15. --- «Ah! Ali está o tio José, que traz um \VUgras{carro de feixes de
lenha} \textsuperscript{(30)} para o padeiro. --- É verdade, disse o meu pai,
reconheço a sua junta de \VUgras{mulas} \textsuperscript{(31)}. Tenho de falar com
ele e vou aproveitar a ocasião. Até à vista, tio Jaime».

%%K t07-c1-16-1 p
16. --- Justamente quando saíamos, as crianças da aldeia deixavam a
escola e lançavam-se a correr para o \VUgras{escorrega} \textsuperscript{(33)},
com o risco de caírem e partirem um membro na \VUgras{queda}
\textsuperscript{(34)}. Uma delas foi buscar o seu \VUgras{trenó}
\textsuperscript{(32)} a casa do carpinteiro de carros, sentou dentro um dos seus
companheiros e partiu a correr.

%%K t07-c1-17-1 p
17. --- Outras precipitaram-se para um \VUgras{monte de neve}
\textsuperscript{(37)} junto à \VUgras{ponte} \textsuperscript{(40)} e começaram a
bombardear o \VUgras{boneco de neve} \textsuperscript{(39)} que tinham feito na
véspera e ao qual tinham posto um chapéu, um cachimbo e uma mochila
escolar a tiracolo.

%%K t07-c1-18-1 p
18. --- Pouco lhes importavam o vento gelado e o frio. A maioria nem
sequer tinha \VUgras{cachecol} \textsuperscript{(35)} e, para se aquecerem depois
da batalha de \VUgras{bolas de neve} \textsuperscript{(38)}, bastava-lhes um
punhado de \VUgras{castanhas} \textsuperscript{(42)} muito quentes compradas à tia
Jacoba, a \VUgras{vendedeira}, que tirita de frio sob o seu \VUgras{xaile}
\textsuperscript{(43)} demasiado fino.

%%K t07-tit-5 sub
\VUsaut{3.0mm}
\VUcentre{12.0pt}{0}{\textit{Segunda cena.}}
\VUsaut{3.0mm}

%%K t07-tit-6 sub
\VUpk{10.2pt}{40}{A casa de lavoura na primavera.}
\VUsaut{4.0mm}

%%K t07-c2-01-1 p
1. --- Apesar de todos os prazeres do inverno, com fundo gozo saudamos a
volta da \VUgras{primavera}.

%%K t07-c2-01-2 p
Que quadro tão alegre é o terreiro de uma quinta, ao entardecer, no mês
de maio!

%%K t07-c2-02-1 p
2. --- No horizonte, o \VUgras{sol} \textsuperscript{(1)} põe-se devagar, inundando
o \VUgras{campo} \textsuperscript{(3)} com os seus \VUgras{raios} \textsuperscript{(2)} de
púrpura. O diligente \VUgras{lavrador} \textsuperscript{(6)} despacha-se a lavrar o
seu \VUgras{campo} \textsuperscript{(4)}.

%%K t07-c2-02-2 p
Com o seu \VUgras{arado} \textsuperscript{(5)}, puxado por um vigoroso
\VUgras{cavalo} \textsuperscript{(7)}, traça o \VUgras{sulco} \textsuperscript{(8)} no qual
o \VUgras{semeador} \textsuperscript{(9)} deitará às mãos-cheias a
\VUgras{semente} que dará as espigas douradas.

%%K t07-c2-03-1 p
3. --- Os \VUgras{ceifeiros} \textsuperscript{(11)}, com as suas gadanhas, cortaram
a erva do prado, os fenadores secaram o \VUgras{feno} \textsuperscript{(10)}
voltando-o com os \VUgras{forcados} \textsuperscript{(12)} e os ancinhos, e ei-los
de volta.

%%K t07-c2-03-2 p
Uns vão diante do pesado carro puxado por bois; levam a sua ferramenta ao
ombro. Outros, mais preguiçosos, instalaram-se comodamente sobre o feno
cheiroso.

%%K t07-c2-04-1 p
4. --- Ao passar saúdam amistosamente o \VUgras{jardineiro}
\textsuperscript{(16)}, ocupado no \VUgras{pomar} \textsuperscript{(13)} a podar as
\VUgras{árvores de fruto} e a enxertar as \VUgras{pereiras}
\textsuperscript{(14)}. As \VUgras{ameixeiras} \textsuperscript{(15)} estão em flor e, na
\VUgras{horta} \textsuperscript{(17)} bem estrumada, os legumes, couves e alfaces
já estão belos.

%%K t07-c2-04-2 p
Sobre o muro baixo, um belo \VUgras{pavão} \textsuperscript{(18)} abre a cauda para
que lhe admirem as magníficas penas.

%%K t07-tit-7 sub
\VUpk{10.2pt}{40}{O terreiro da quinta.}
\VUsaut{3.0mm}

%%K t07-c2-05-1 p
5. --- As portas da \VUgras{cavalariça} \textsuperscript{(19)} estão abertas e
veem-se a manjedoura e a \VUgras{cama de palha} \textsuperscript{(23)} da
\VUgras{égua} \textsuperscript{(21)} que o \VUgras{moço de cavalariça}
\textsuperscript{(20)} acaba de desatrelar. O \VUgras{poldro} \textsuperscript{(22)}, tão
gracioso com as suas pernas compridas, segue a mãe aos saltos.

%%K t07-c2-06-1 p
6. --- Junto ao \VUgras{poço} \textsuperscript{(29)}, as \VUgras{vacas}
\textsuperscript{(25)} vão beber ao \VUgras{gamelão} \textsuperscript{(26)} de madeira que
lhes serve de bebedouro. O \VUgras{vitelo} \textsuperscript{(24)} aproveita para
mamar, e o \VUgras{boi} \textsuperscript{(27)}, sentindo o bom cheiro da forragem,
muge alegremente e levanta com orgulho a cabeça de poderosos
\VUgras{cornos} \textsuperscript{(28)}.

%%K t07-c2-07-1 p
7. --- Todos trabalham. A \VUgras{criada} \textsuperscript{(32)} segura com a mão a
\VUgras{corda} \textsuperscript{(31)} do poço. Tira água com um \VUgras{balde}
\textsuperscript{(30)} para dar de beber ao gado. Junto ao \VUgras{celeiro}
\textsuperscript{(33)}, um homem ajunta com o ancinho as palhinhas, e diante da
porta da \VUgras{casa} \textsuperscript{(34)}, uma velha \VUgras{fiandeira}
\textsuperscript{(35)}, com a sua roca e a sua \VUgras{estriga}, fia o \VUgras{linho}
ou o \VUgras{cânhamo} como noutros tempos.

%%K t07-tit-8 sub
\VUsaut{3.0mm}
\VUpk{10.2pt}{40}{O lavrador.}
\VUsaut{3.0mm}

%%K t07-c2-08-1 p
8. --- O \VUgras{lavrador} \textsuperscript{(36)} dirige todo este trabalho e dá
ordens aos seus criados. Manda recolher a \VUgras{grade} \textsuperscript{(40)} e a
\VUgras{picareta} \textsuperscript{(39)} que se esqueceram junto à \VUgras{casota}
\textsuperscript{(38)} do \VUgras{cão} \textsuperscript{(37)}.

%%K t07-c2-09-1 p
9. --- Os seus gritos assustam uma \VUgras{pomba} \textsuperscript{(41)}, que foge a
toda a pressa para o \VUgras{pombal} \textsuperscript{(42)}, e as \VUgras{galinhas}
\textsuperscript{(43)}, que se apressam a voltar ao \VUgras{galinheiro}
\textsuperscript{(44)}.

%%K t07-c2-10-1 p
10. --- Diante do \VUgras{redil} \textsuperscript{(48)}, eis o velho \VUgras{pastor}
\textsuperscript{(45)} que traz de volta o seu \VUgras{rebanho} \textsuperscript{(49)}.
Com o seu \VUgras{cajado} \textsuperscript{(46)}, que nunca o deixa, tal como o seu
\VUgras{blusão} \textsuperscript{(47)} descolorido, leva ao aprisco
\VUgras{carneiros} \textsuperscript{(50)}, \VUgras{ovelhas} \textsuperscript{(53)},
\VUgras{cordeiros} \textsuperscript{(54)}, \VUgras{cabras} \textsuperscript{(51)} e
\VUgras{cabritos} \textsuperscript{(52)}. O seu fiel \VUgras{cão}
\textsuperscript{(55)} Argos fecha a marcha e anima com os seus latidos os
retardatários.

%%K t07-c2-11-1 p
11. --- Todo este ruído e este movimento não perturbam a calma da boa
\VUgras{vaca leiteira} \textsuperscript{(56)}, que faz soar o seu \VUgras{chocalho}
\textsuperscript{(57)} enquanto a ordenham diante da porta do \VUgras{estábulo}
\textsuperscript{(58)}.

%%K t07-c2-12-1 p
12. --- Parece compreender que tem de dar o seu leite para que a
\VUgras{lavradeira} \textsuperscript{(60)} possa fazer \VUgras{manteiga} na
\VUgras{batedeira} \textsuperscript{(61)} ou os excelentes \VUgras{queijos}
\textsuperscript{(64)} que escorrem sobre a tábua posta em cima da
\VUgras{gamela} \textsuperscript{(63)}.

%%K t07-c2-13-1 p
13. --- Toda a \VUgras{leitaria} \textsuperscript{(59)} é mantida muito limpa pelo
\VUgras{moço da quinta} \textsuperscript{(65)}, que acaba de lavar os
\VUgras{cântaros do leite} \textsuperscript{(62)} no grande balde que leva na mão.

%%K t07-tit-9 sub
\VUsaut{3.0mm}
\VUpk{10.2pt}{40}{A capoeira.}
\VUsaut{3.0mm}

%%K t07-c2-14-1 p
14. --- Com os seus grossos tamancos anda com alguma falta de jeito, e
assim faz fugir o \VUgras{peru} \textsuperscript{(68)}, as \VUgras{patas}
\textsuperscript{(66)}, os \VUgras{patos} \textsuperscript{(67)} e os \VUgras{gansos}
\textsuperscript{(69)}, que abrem muito o \VUgras{bico} \textsuperscript{(70)} e lançam
gritos inquietos.

%%K t07-c2-14-2 p
O \VUgras{galo} \textsuperscript{(71)}, mais belicoso, levanta a sua \VUgras{crista}
\textsuperscript{(72)} vermelha, sacode as penas da sua \VUgras{cauda}
\textsuperscript{(73)} e lança um sonoro «cocorocó».

%%K t07-c2-15-1 p
15. --- Uma \VUgras{galinha mãe} \textsuperscript{(74)} debica no \VUgras{monturo}
\textsuperscript{(81)} com os seus \VUgras{pintainhos} \textsuperscript{(75)}. O
\VUgras{leitão} \textsuperscript{(78)} corre para a mãe, a enorme \VUgras{porca}
\textsuperscript{(77)} que vemos junto à pocilga.

%%K t07-c2-16-1 p
16. --- Nas suas gaiolas, os \VUgras{coelhos} \textsuperscript{(79)} comem com
avidez a \VUgras{erva} \textsuperscript{(80)} tenra que lhes traz a filha do
lavrador. Joaninha gosta muito dos seus coelhos e, todos os dias, depois
de ir buscar os \VUgras{ovos} \textsuperscript{(76)} frescos ao galinheiro, vem
visitar os seus amiguinhos.
\VUsaut{6.0mm}
\VUfilet{22mm}
